\section{Sound by construction}

The previous chapter's pipeline (build the net, derive $N^*$, check liveness and boundedness) is \emph{verification}: design first, check later. The opposite is \emph{construction}: assemble the workflow from pre-validated pieces under composition rules that preserve soundness, so the result is sound \emph{by construction}. Two ingredients: a catalogue of \textbf{building blocks}, small WF nets easily proved sound and \emph{safe} (1-bounded), and \textbf{composition patterns} under which sound-and-safe is closed.

\subsection{Composition by refinement}

Let $N$ and $N'$ be safe and sound WF nets with interfaces $i, o$ and $i', o'$, and let $t$ be a task of $N$ with \textbf{exactly one input and one output place} ($|{}^\bullet t| = |t^\bullet| = 1$).

\begin{definition}[{Refinement $N[N'/t]$}]
$N[N'/t]$ is obtained by \emph{replacing $t$ with $N'$}: the unique input place of $t$ is fused with $i'$, the unique output place with $o'$, and $t$ with its two arcs is removed: macro-expansion of a task into a full sub-procedure.
\end{definition}

\begin{theorem}[Compositionality]
If $N$ and $N'$ are sound and safe, so is $N[N'/t]$. (Proof sketched below.)
\end{theorem}

\begin{definitionbox}{Initial and final transitions}
A transition $t$ is \emph{initial} if $i \in {}^\bullet t$ and \emph{final} if $o \in t^\bullet$. In a \emph{sound} WF net:
\begin{itemize}
  \item every initial transition has ${}^\bullet t = \{i\}$: were there another input place $p$, no token could ever reach $p$ before $t$ fires, so $t$ would be dead;
  \item every final transition has $t^\bullet = \{o\}$: were there another output place, its token would either survive completion (violating proper completion) or never be consumable.
\end{itemize}
\end{definitionbox}

\emph{Proof idea of the theorem}: a sound WF net behaves like a transition whose consumption and production are not atomic: one token in, eventually one token out, nothing else across the interface. Formally one exhibits a bijection between the reachable markings of $N[N'/t]$ and the interface-respecting pairs of reachable markings of $N$ and $N'$; safeness and soundness of the composite follow from those of the parts.

\subsection{Building blocks}

\begin{center}
\begin{tikzpicture}[node distance=0.25cm and 0.4cm]
  \node[bpmn event] (b1) {};
  \node[bpmn task, rounded corners=0pt, right=of b1, minimum width=0.5cm, minimum height=0.4cm] (bt) {\scriptsize $t$};
  \node[bpmn event, right=of bt] (b2) {};
  \node[right=0.2cm of b2, font=\scriptsize] {basic};
  \draw[bpmn flow] (b1) -- (bt);
  \draw[bpmn flow] (bt) -- (b2);
  \begin{scope}[xshift=4.4cm]
  \node[bpmn event] (s1) {};
  \node[bpmn task, rounded corners=0pt, right=of s1, minimum width=0.5cm, minimum height=0.4cm] (st) {\scriptsize $t$};
  \node[bpmn event, right=of st] (s2) {};
  \node[bpmn task, rounded corners=0pt, right=of s2, minimum width=0.5cm, minimum height=0.4cm] (st2) {\scriptsize $t'$};
  \node[bpmn event, right=of st2] (s3) {};
  \node[right=0.2cm of s3, font=\scriptsize] {sequence};
  \draw[bpmn flow] (s1) -- (st);
  \draw[bpmn flow] (st) -- (s2);
  \draw[bpmn flow] (s2) -- (st2);
  \draw[bpmn flow] (st2) -- (s3);
  \end{scope}
\end{tikzpicture}\\[6pt]
\begin{tikzpicture}[node distance=0.25cm and 0.4cm]
  \node[bpmn event] (x1) {};
  \node[bpmn task, rounded corners=0pt, above right=0.1cm and 0.5cm of x1, minimum width=0.5cm, minimum height=0.4cm] (xt) {\scriptsize $t$};
  \node[bpmn task, rounded corners=0pt, below right=0.1cm and 0.5cm of x1, minimum width=0.5cm, minimum height=0.4cm] (xt2) {\scriptsize $t'$};
  \node[bpmn event, right=2.0cm of x1] (x2) {};
  \node[right=0.2cm of x2, font=\scriptsize] {implicit XOR};
  \draw[bpmn flow] (x1) -- (xt);
  \draw[bpmn flow] (x1) -- (xt2);
  \draw[bpmn flow] (xt) -- (x2);
  \draw[bpmn flow] (xt2) -- (x2);
  \begin{scope}[xshift=6cm]
  \node[bpmn event] (i1) {};
  \node[bpmn task, rounded corners=0pt, right=of i1, minimum width=0.5cm, minimum height=0.4cm] (it) {\scriptsize $t$};
  \node[bpmn event, right=of it] (i2) {};
  \node[bpmn task, rounded corners=0pt, below=0.3cm of it, minimum width=0.5cm, minimum height=0.4cm] (it2) {\scriptsize $t'$};
  \node[right=0.2cm of i2, font=\scriptsize] {iteration};
  \draw[bpmn flow] (i1) -- (it);
  \draw[bpmn flow] (it) -- (i2);
  \draw[bpmn flow] (i2) to[bend left=20] (it2);
  \draw[bpmn flow] (it2) to[bend left=20] (i1);
  \end{scope}
\end{tikzpicture}
\end{center}
All four are trivially sound and safe. The \emph{implicit XOR} branches on a free-choice place (the token enables both alternatives, the chosen one consumes it); the two internal places of the \emph{iteration} block must not coincide with $i$ or $o$ of the host net.

A second family reifies branching as transitions. Here one caution: a single transition never makes a choice: firing it marks \emph{all} its output places at once, which is an AND, not an XOR. The choice must sit on a place. So the \emph{explicit XOR-split} is an entry transition feeding one internal place from which two alternative transitions compete (the choice again sits on the place); the \emph{explicit XOR-join} is the dual; and the combined \emph{XOR block} chains an explicit split, two alternative tasks and an explicit merge. The \emph{AND block} is drawn like the XOR block but its entry transition produces \emph{both} internal places simultaneously: $t$ and $t'$ then interleave freely and the exit transition consumes both. The catalogue is open: any net proved sound and safe may join it.

% Verified 2026-07-10: AND block drawn. Places: input i, two internal (post-entry)
% feeding t and t', two internal (post-task) feeding the exit, output o (6 places).
% Transitions: entry (unlabeled), t, t', exit (unlabeled) (4 transitions). Arcs:
% i->entry, entry->{p_t,p_t'}, p_t->t, p_t'->t', t->{q_t}, t'->{q_t'}, q_t->exit,
% q_t'->exit, exit->o. Replayed from M0={i}: fire entry -> {p_t,p_t'}; t and t'
% concurrently enabled, fire t -> {q_t,p_t'}, fire t' -> {q_t,q_t'} (either order);
% exit enabled -> fire -> {o}. Safe (no place ever exceeds 1 token) and sound
% (unique final marking {o}, proper completion). Entry post-set = both internal
% places => concurrency (AND), not choice. Places/non-task transitions left
% unlabeled, matching the implicit-XOR and AND-with-sync figures. Fits A4, no resizebox.
\begin{center}
\begin{tikzpicture}[node distance=0.3cm and 0.55cm]
  \node[bpmn event] (i) {};
  \node[bpmn task, rounded corners=0pt, right=of i, minimum width=0.5cm, minimum height=0.7cm] (en) {};
  \node[bpmn event, above right=0.25cm and 0.5cm of en] (pt) {};
  \node[bpmn event, below right=0.25cm and 0.5cm of en] (pu) {};
  \node[bpmn task, rounded corners=0pt, right=0.5cm of pt, minimum width=0.5cm, minimum height=0.4cm] (t) {\scriptsize $t$};
  \node[bpmn task, rounded corners=0pt, right=0.5cm of pu, minimum width=0.5cm, minimum height=0.4cm] (u) {\scriptsize $t'$};
  \node[bpmn event, right=0.5cm of t] (qt) {};
  \node[bpmn event, right=0.5cm of u] (qu) {};
  \node[bpmn task, rounded corners=0pt, right=0.5cm of qt, yshift=-0.45cm, minimum width=0.5cm, minimum height=0.7cm] (ex) {};
  \node[bpmn event, right=of ex] (o) {};
  \draw[bpmn flow] (i) -- (en);
  \draw[bpmn flow] (en) -- (pt); \draw[bpmn flow] (en) -- (pu);
  \draw[bpmn flow] (pt) -- (t); \draw[bpmn flow] (pu) -- (u);
  \draw[bpmn flow] (t) -- (qt); \draw[bpmn flow] (u) -- (qu);
  \draw[bpmn flow] (qt) -- (ex); \draw[bpmn flow] (qu) -- (ex);
  \draw[bpmn flow] (ex) -- (o);
\end{tikzpicture}
\end{center}

\subsection{Refinement and abstraction at work}

\emph{Refinement} grows nets top-down. Start from the basic block and repeatedly replace a task by a sequence, iteration or XOR block. After a few rounds the net is a structured mix of sequences, a loop and an exclusive branch: sound and safe at every step, with no global check ever run.

\emph{Abstraction} is the dual, and the practical proof technique: to certify a given net, recognise sub-patterns matching catalogue blocks, collapse each to a single task, repeat. When the last collapse reaches the \emph{basic} block, that reduction is the net's structural certificate of soundness and safeness.
\par\smallskip
The word \emph{abstraction} carries three senses across the course, and this calculus is only the first. \textbf{Horizontal} abstraction moves \emph{along} the control flow at one level, folding a block into a single node: the inverse of refinement, exactly the reduction above. \textbf{Vertical} abstraction moves \emph{across} hierarchy levels, hiding a whole sub-net behind one task (the sub-processes of the Petri-net and BPMN chapters). \textbf{Aggregated} abstraction folds \emph{many elements into one} summary (the KPI functions of quantitative analysis, a coarser-grained mined map). The first two are reversible, collapse then expand recovers the original; aggregation is lossy, a summary cannot be resolved back into what it condensed.

\begin{examplebox}{Car Damage claim}
The car-damage claim workflow (drawn as a workflow net in the Petri-nets chapter) was built deliberately by composition, and the abstraction calculus certifies it: the \emph{simple} branch contains a textbook AND block (AND-split, \emph{phone garage} $\parallel$ \emph{check insurance}, AND-join); the \emph{complex} branch is a plain sequence (\emph{check insurance ok}, \emph{check history}, \emph{phone garage}); \emph{classify} with the two branches and their merge forms an explicit XOR; and \emph{decide} routing to \emph{pay} or \emph{send letter}, merged at \emph{end}, forms a second one. Every pattern is in the catalogue, the net reduces to the basic block: \textbf{sound and safe by construction}.
\end{examplebox}

% Rewritten: the previous version claimed the exercise net leaves an
% ``asymmetric AND residue'' where ``the calculus is stuck''. It does not.
% Slides p.35--50 reduce the net all the way to the basic block; slides
% p.51--54 show a second net that, once desugared, is the same net and
% reduces too. The real lesson is the opposite of a stall: the method
% extends via desugaring. Corrected to match the slides.
\begin{examplebox}{Desugaring first, then reduce}
A more intricate exercise net does not match any catalogue block as drawn: its XOR and AND branch points are still gateway shorthands, not explicit transitions. The fix is \emph{desugaring}: expand each gateway into the Petri-net transitions it abbreviates (an XOR-split into two alternative entry transitions on a shared place, an AND-split into one forking transition, and so on). Once desugared, the net is a plain Petri net of sequences, XOR blocks, an AND block and an iteration, and abstraction certifies it clause by clause: a sequence collapses, an implicit XOR collapses, the lower-branch iteration collapses, the two parallel branches fold into the AND block, and the last step reaches the \emph{basic} block: \textbf{sound and safe}. A second exercise looks different but, after the same desugaring, \emph{is the same net}, and reduces identically. The moral is not that the calculus stalls: it is that syntax must be normalised before the block patterns become visible. Desugaring is the method's first degree of freedom; the second (extending the catalogue with independently verified blocks) is the generalisation below.
\end{examplebox}

\subsection{Generalised substitution}

The substitution rule so far requires $|{}^\bullet t| = |t^\bullet| = 1$. The generalisation targets transitions with arbitrary fan-in and fan-out:
\begin{center}
\begin{tikzpicture}[node distance=0.25cm and 0.5cm]
  \node[bpmn event] (p1) {};
  \node[bpmn event, below=0.4cm of p1] (p2) {};
  \node[bpmn task, rounded corners=0pt, right=0.7cm of p1, yshift=-0.45cm, minimum width=0.55cm, minimum height=0.45cm] (t) {\scriptsize $t$};
  \node[bpmn event, right=0.7cm of t, yshift=0.45cm] (q1) {};
  \node[bpmn event, below=0.4cm of q1] (q2) {};
  \draw[bpmn flow] (p1) -- (t);
  \draw[bpmn flow] (p2) -- (t);
  \draw[bpmn flow] (t) -- (q1);
  \draw[bpmn flow] (t) -- (q2);
\end{tikzpicture}
\end{center}

\begin{definitionbox}{General replacement $N[N'/t]$}
Let $t$ be any transition of $N$, and $N'$ a WF net with initial place $i'$, final place $o'$,
\[
  T_{i'} = \{u \mid {}^\bullet u = \{i'\}\} \quad \text{(initial transitions)}, \qquad
  T_{o'} = \{v \mid v^\bullet = \{o'\}\} \quad \text{(final transitions)}.
\]
$N[N'/t]$ removes $t$ from $N$ and $i', o'$ from $N'$, then wires: for every input arc $(p, t)$ of $t$ and every $u \in T_{i'}$, an arc $(p, u)$; for every output arc $(t, q)$ and every $v \in T_{o'}$, an arc $(v, q)$.
\par\smallskip
\textbf{Theorem} (proof omitted). If $N$ and $N'$ are sound and safe, so is $N[N'/t]$.
\end{definitionbox}

The single-input single-output rule is the special case where ${}^\bullet t$ and $t^\bullet$ are singletons. A fifth catalogue block exploits the new freedom, \emph{AND with sync}: an AND-split forks two places, a central transition consumes \emph{both} and produces \emph{both} the places consumed by the final AND-join:
\begin{center}
\begin{tikzpicture}[node distance=0.3cm and 0.55cm]
  \node[bpmn event] (i) {};
  \node[bpmn task, rounded corners=0pt, right=of i, minimum width=0.5cm, minimum height=0.4cm] (sp) {};
  \node[bpmn event, above right=0.25cm and 0.5cm of sp] (a1) {};
  \node[bpmn event, below right=0.25cm and 0.5cm of sp] (a2) {};
  \node[bpmn task, rounded corners=0pt, right=2.0cm of sp, minimum width=0.5cm, minimum height=0.7cm] (sync) {\scriptsize $t$};
  \node[bpmn event, above right=0.25cm and 0.5cm of sync] (b1) {};
  \node[bpmn event, below right=0.25cm and 0.5cm of sync] (b2) {};
  \node[bpmn task, rounded corners=0pt, right=2.0cm of sync, minimum width=0.5cm, minimum height=0.4cm] (jn) {};
  \node[bpmn event, right=of jn] (o) {};
  \draw[bpmn flow] (i) -- (sp);
  \draw[bpmn flow] (sp) -- (a1); \draw[bpmn flow] (sp) -- (a2);
  \draw[bpmn flow] (a1) -- (sync); \draw[bpmn flow] (a2) -- (sync);
  \draw[bpmn flow] (sync) -- (b1); \draw[bpmn flow] (sync) -- (b2);
  \draw[bpmn flow] (b1) -- (jn); \draw[bpmn flow] (b2) -- (jn);
  \draw[bpmn flow] (jn) -- (o);
\end{tikzpicture}
\end{center}

\begin{examplebox}{Generalised refinement and a full certification}
\emph{Refinement with fan-in/fan-out}: a central transition $t_2$ with two inputs ($p_2, p_3$) and two outputs ($p_4, p_5$) is replaced by a sequence $u \to q \to v$; being the unique initial and final transitions of the inserted net, $u$ inherits both input arcs and $v$ both output arcs: the interface is preserved, the behaviour expands.
\par\smallskip
\emph{A complete reduction}: a larger multi-arc net is certified by abstraction: an implicit XOR in the upper branch collapses, then two sequences; the lower branch contains a non-canonical sub-net (two parallel multi-arc transitions plus a feedback transition) that is promoted to an \emph{ad-hoc building block} after an independent proof of its soundness and safeness; with it in the catalogue the lower branch reduces to a sequence, then to a single transition; the two residual branches between $i$ and $o$ form an implicit XOR; one last collapse reaches the basic block. Both degrees of freedom are at work here, taking the calculus beyond the four primitive patterns.
\end{examplebox}
